\documentclass[instructions.tex]{subfiles}
\begin{document}
 
\chapter{Les occasions de l’impureté et ses remèdes.}
 
\primo Toute personne qui pense sérieusement aux suites de l’impureté, dit saint Martin de Braga, tâche d’en éviter l’occasion et les plus légères atteintes, qui devant Dieu sont d’une plus grande conséquence qu’on ne le pense. Ce qu’on regarde comme une bagatelle en cette matière, est quelquefois le commencement de la damnation\footnote{Libidinis initia continebit, qui exitum cogitabit.}. Les occasions et les pièges les plus ordinaires sont un enjouement affecté dans les parures, la liberté des sens, les fréquentations.
 
1. L’enjouement et la parure. Jeune homme, qui vous parez pour inspirer de la passion aux personnes du sexe, qui n’avez avec elles que des manières flatteuses et des airs complaisans, vous êtes leur tentateur et leur démon, dit saint Clément. Elles ne sont pas moins faibles que vous, et si leur cœur est souillé, vous êtes le meurtrier de leurs âmes. Et vous, filles et femmes, malheur à vous, si vous affectez d’attirer sur vous les regards d’autrui ! Plus vous voulez paraître agréables au monde, plus vous êtes abominables devant Dieu. Une personne du sexe enjouée, vêtue sans modestie, est l’organe du démon, dit saint Bernard; c’est par elle qu’il tente, qu’il parle\footnote{Organum Satanæ. S. Bern.}. Elle est, dit saint Cyprien, comme une forteresse où le démon est en embuscade pour surprendre les âmes.
 
2. La liberté des sens. Veillez sur vous, et veillez toujours, dit Jésus-Christ. La mort entre par les fenêtres, c’est-à-dire que l’impureté, le poison de l’âme, entre dans le cœur par les yeux, par les oreilles, par les paroles, par les chansons, par les portraits, par les lectures. Il y entre par les comédies, par les danses et les spectacles; c’est là que, l’esprit dissipé et le cœur ému, l’âme goûte le poison impur sans y prendre garde. Plus ce poison paraît doux, plus il est subtil et mortel. Il tue souvent l’âme aussitôt qu’on le regarde et qu’on en approche. Un regard, une pensée, un désir, la souillent, et sont capables de vous perdre. La vue d’une femme a souvent vaincu ceux que les tentations et les persécutions n’avaient pu vaincre.
 
Les fréquentations. Qu’aucun homme, jeune ou âgé, écrivait saint Bernard à sa sœur, n’ait avec vous aucune société, aucune assiduité, aucune familiarité, quelque saint et quelque régulier qu’il soit. Les assiduités et certaines promenades avec des personnes d’un sexe différent sont toujours dangereuses, rarement innocentes. Les libertés familières, les embrassemens folâtres et les cajoleries sont les arrhes de l’impureté et les marques d’une chasteté mourante ou qui est déjà morte, dit saint Jérôme. Que les jeunes gens sont à plaindre, s’ils ne connaissent pas leurs dangers! et que les parens sont coupables, s’ils ne veillent pas à leur sûreté !
 
\secundo Les remèdes pour se conserver chaste sont, 1. la résolution et la crainte de Dieu. Une courtisane étant entrée dans la chambre de saint Thomas d’Aquin pour le séduire, ce saint jeune homme prit un tison, et la mit en fuite, et depuis il n’eut plus aucune tentation à ce sujet. Ayez du courage, recourez à la prière, craignez Dieu dans les occasions, et vous vous procurerez des grâces pour vous soutenir.
 
2. La fidélité dans les tentations. Si vous avez de fréquentes tentations, ne vous en étonnez pas. Les saints en avaient ainsi que vous; c’est leur fidélité à y résister qui les a rendus saints. Les pensées que vous avez malgré vous, loin de vous nuire, sont un sujet de mérite et de gloire. Aussitôt que la pensée se présente à votre esprit, donnez-lui le change, pensez promptement à d’autres choses. Ne balancez point: priez, gémissez, jetez-vous en esprit au pied du crucifix, implorez le secours de la mère de Dieu, de votre saint ange et de vos saints protecteurs.
 
3. Le travail et la mortification. Prenez garde de donner à votre corps, par l’oisiveté, par la mollesse et l’intempérance, des armes pour vous faire la guerre. Votre corps est un esclave qu’il faut traiter durement. Apprenez, en voyant le corps innocent de Jésus, couvert de sang, comment vous devez traiter le vôtre qui est coupable. Ne vous fiez ni à votre vertu, ni à votre âge. Salomon, le plus sage des hommes, fut séduit par les femmes dans sa vieillesse.
 
4. Adressez-vous à un confesseur qui vous fasse quitter vos connaissances dangereuses, vos habitudes, vos attaches, et qui vous dispose à recevoir dignement la sainte Eucharistie, que l’Ecriture appelle le froment des élus, et le vin qui fait les vierges. Combien de gens maudissent en enfer les confesseurs qui les ont laissés vivre dans le sacrilége, dans le libertinage et la mollesse!
 
Si tout ce qui a été dit ci-devant ne vous touche pas, vous n’avez plus la foi d’un chrétien ni la raison d’un homme; car l’impureté fait perdre la foi, rend l’homme insensé, et fait apostasier les sages. Vous êtes, selon la parole de l’Ecriture, un insensé qui, devenu la victime d’une passion qui vous entraîne, êtes tellement abruti, que vous n’avez plus que les sensations de l’animal\footnote{Tam sequitur quasi bos ductus ad victimam, et quasi agnus lasciviens et ignorans quod ad vincula stultus trahatur. Prov. 7. 22.}.
 
\section{Résolutions}
 
\primo Je m’habillerai toujours d’une manière décente, et je conserverai partout cet extérieur de retenue et de modestie qui impose le respect.
 
\secundo Je veillerai sur mes yeux, sur mes oreilles, sur ma langue, sur tout moi-même, afin de ne voir, ni entendre, ni dire, ni faire de propos délibéré, rien qui puisse blesser la plus délicate des vertus.
 
\tertio Et je n’oublierai jamais que la victoire sur l'ennemi de la chasteté dépend en grande partie de la diligence avec laquelle on lui resiste dès le premier instant du combat. 

\prayer{Mon Dieu, fortifiez-moi dans mes résolutions; et faites-moi la grâce d'éviter, avec tout le soin possible, un péché qui vous déplait si horriblement.}
 
\end{document}
